\documentclass[aps,preprint]{revtex4}%
\usepackage{amssymb}
\usepackage{amsfonts}
\usepackage{amsmath}
\usepackage{graphicx}%
\setcounter{MaxMatrixCols}{30}
%TCIDATA{OutputFilter=latex2.dll}
%TCIDATA{Version=5.50.0.2960}
%TCIDATA{CSTFile=revtex4.cst}
%TCIDATA{Created=Tuesday, June 12, 2018 08:30:22}
%TCIDATA{LastRevised=Friday, May 03, 2019 07:55:05}
%TCIDATA{<META NAME="GraphicsSave" CONTENT="32">}
%TCIDATA{<META NAME="SaveForMode" CONTENT="1">}
%TCIDATA{BibliographyScheme=Manual}
%TCIDATA{<META NAME="DocumentShell" CONTENT="Articles\SW\REVTeX 4">}
%BeginMSIPreambleData
\providecommand{\U}[1]{\protect\rule{.1in}{.1in}}
%EndMSIPreambleData
\newtheorem{theorem}{Theorem}
\newtheorem{acknowledgement}[theorem]{Acknowledgement}
\newtheorem{algorithm}[theorem]{Algorithm}
\newtheorem{axiom}[theorem]{Axiom}
\newtheorem{claim}[theorem]{Claim}
\newtheorem{conclusion}[theorem]{Conclusion}
\newtheorem{condition}[theorem]{Condition}
\newtheorem{conjecture}[theorem]{Conjecture}
\newtheorem{corollary}[theorem]{Corollary}
\newtheorem{criterion}[theorem]{Criterion}
\newtheorem{definition}[theorem]{Definition}
\newtheorem{example}[theorem]{Example}
\newtheorem{exercise}[theorem]{Exercise}
\newtheorem{lemma}[theorem]{Lemma}
\newtheorem{notation}[theorem]{Notation}
\newtheorem{problem}[theorem]{Problem}
\newtheorem{proposition}[theorem]{Proposition}
\newtheorem{remark}[theorem]{Remark}
\newtheorem{solution}[theorem]{Solution}
\newtheorem{summary}[theorem]{Summary}
\newenvironment{proof}[1][Proof]{\noindent\textbf{#1.} }{\ \rule{0.5em}{0.5em}}
\begin{document}
\preprint{HEP/123-qed}
\title[Short title for running header]{REV\TeX4}
\author{First Author}
\affiliation{My Institution}
\author{Second Author}
\affiliation{My Institution}
\author{Third Author}
\affiliation{Other Institution}
\keywords{one two three}
\pacs{PACS number}

\begin{abstract}
Shell document for REV\TeX{} 4.

\end{abstract}
\volumeyear{year}
\volumenumber{number}
\issuenumber{number}
\eid{identifier}
\date[Date text]{date}
\received[Received text]{date}

\revised[Revised text]{date}

\accepted[Accepted text]{date}

\published[Published text]{date}

\startpage{101}
\endpage{102}
\maketitle
\tableofcontents


\section{Clifford Algebra}

A Clifford map $\Phi$ from an \emph{inner product} vector space $E$ to an
associative algebra $\mathcal{A}$
\begin{equation}
\Phi:E\rightarrow\mathcal{A}%
\end{equation}
is given by
\begin{equation}
\Phi\left(  x\right)  ^{2}=\left(  \left.  x\right\vert x\right)
I_{\mathcal{A}},
\end{equation}
which is equivalent to
\begin{equation}
\Phi\left(  x\right)  \Phi\left(  y\right)  +\Phi\left(  y\right)  \Phi\left(
x\right)  =2\left(  \left.  x\right\vert y\right)  I_{\mathcal{A}}.
\end{equation}
The subspace $\Phi\left(  E\right)  $ generates the Clifford algebra
$\mathcal{C}\left(  E\right)  \subseteq\mathcal{A}$ and as the map $\Phi$ is
linar and injective one can identify $E$ and $\Phi\left(  E\right)  $ and set
\begin{equation}
xy+yx=2\left(  \left.  x\right\vert y\right)  I_{\mathcal{A}},\,x,y\in E.
\end{equation}


\subsubsection{Examples}

\begin{enumerate}
\item $E=\mathbb{R};\left(  \left.  x\right\vert y\right)  =-xy$%
\begin{align}
\Phi &  :\mathbb{R}\rightarrow\mathbb{C}\\
\mathcal{C}\left(  E\right)   &  =\mathbb{C}.
\end{align}


\item $E=\mathbb{R}^{2};\left(  \left.  x\right\vert y\right)  =-\sum
\limits_{1\leq j\leq2}x_{j}y_{j}$%
\begin{align}
\Phi &  :\mathbb{R}^{2}\rightarrow\mathbb{H}\\
\mathcal{C}\left(  E\right)   &  =\mathbb{H}.
\end{align}


\item dim$\left\{  E\right\}  =r;\left(  \left.  x\right\vert y\right)
=0;\forall x,y$%
\begin{equation}
\mathcal{C}\left(  E\right)  =\wedge\left(  E\right)  .
\end{equation}

\end{enumerate}

\subsubsection{$\mathbb{Z}_{2}$-Grading of Clifford Algebra}

One can decompose the Clifford algebra $\mathcal{C}\left(  E\right)  $ into a
direct sum of a subalgebra, $\mathcal{C}\left(  E\right)  _{0},$ and a
subspace $\mathcal{C}\left(  E\right)  _{1}$ with a $\mathbb{Z}_{2}$ grading
\begin{equation}
\mathcal{C}\left(  E\right)  =\mathcal{C}\left(  E\right)  _{even}%
\oplus\mathcal{C}\left(  E\right)  _{odd}\equiv\mathcal{C}\left(  E\right)
_{0}\oplus\mathcal{C}\left(  E\right)  _{1},
\end{equation}
associated with an involution
\begin{equation}
\omega\left(  x\right)  =-x,
\end{equation}
that extends to all of $\mathcal{C}\left(  E\right)  $ by the definition
\begin{equation}
\omega\left(  x_{1}\cdots x_{n}\right)  =\omega\left(  x_{1}\right)
\cdots\omega\left(  x_{n}\right)  ,
\end{equation}
to give
\begin{align}
\omega &  :\mathcal{C}\left(  E\right)  _{0}\rightarrow\mathcal{C}\left(
E\right)  _{0}\\
\omega &  :\mathcal{C}\left(  E\right)  _{1}\rightarrow\mathcal{C}\left(
E\right)  _{1}.
\end{align}


Consider two inner product spaces $E$ and $F$ and their direct sum $E$
$\oplus$ $F:$

\begin{theorem}
The $\mathbb{Z}_{2}$-graded algebra $\mathcal{C}\left(  E\oplus F\right)  $ is
isomorphic to the skew tensor product of the $\mathbb{Z}_{2}$-graded algebras
$\mathcal{C}\left(  E\right)  $ and $\mathcal{C}\left(  F\right)  $, i.e.
\begin{equation}
\mathcal{C}\left(  E\oplus F\right)  \cong\mathcal{C}\left(  E\right)
\widehat{\otimes}\mathcal{C}\left(  F\right)  .
\end{equation}

\end{theorem}

\begin{proof}
Let $H=E\oplus F$ then the injections
\begin{align}
i  &  :E\rightarrow H\\
j  &  :F\rightarrow H
\end{align}
induce%
\begin{align}
i_{C}  &  :\mathcal{C}\left(  E\right)  \rightarrow\mathcal{C}\left(  H\right)
\\
j_{C}  &  :\mathcal{C}\left(  F\right)  \rightarrow\mathcal{C}\left(
H\right)  .
\end{align}
Define a linear map $f$
\begin{equation}
f:\mathcal{C}\left(  E\right)  \widehat{\otimes}\mathcal{C}\left(  F\right)
\rightarrow\mathcal{C}\left(  H\right)
\end{equation}
by
\begin{equation}
f\left(  a\otimes b\right)  =i_{C}\left(  a\right)  \bullet j_{C}\left(
b\right)  ,
\end{equation}
where "$\bullet"$ is Clifford multiplication. The skew tensor product
$\mathcal{C}\left(  E\right)  \widehat{\otimes}\mathcal{C}\left(  F\right)  $
\emph{is contained} in the tensor product $\mathcal{C}\left(  E\right)
\otimes\mathcal{C}\left(  F\right)  ,$ hence
\begin{equation}
f:\mathcal{C}\left(  E\right)  \otimes\mathcal{C}\left(  F\right)
\rightarrow\mathcal{C}\left(  E\right)  \widehat{\otimes}\mathcal{C}\left(
F\right)  .
\end{equation}
One can show that $f$ is an algebra homomorphism and
\begin{equation}
i_{C}\left(  a\right)  \bullet j_{C}\left(  b\right)  =\left(  -1\right)
^{\alpha\beta}i_{C}\left(  b\right)  \bullet j_{C}\left(  a\right)
;\,\alpha=\deg\left\{  a\right\}  ,\beta=\deg\left\{  b\right\}
,a\in\mathcal{C}\left(  E\right)  ,b\in\mathcal{C}\left(  F\right)
\end{equation}
Since
\begin{equation}
i_{C}\left(  E\right)  \bot i_{C}\left(  F\right)
\end{equation}
the commutation relationship gives
\begin{equation}
i_{C}\left(  x\right)  \bullet j_{C}\left(  y\right)  +i_{C}\left(  y\right)
\bullet j_{C}\left(  x\right)  =0\,\forall x\in i_{C}\left(  E\right)
,\forall y\in i_{C}\left(  F\right)
\end{equation}
thus
\begin{align}
f\left(  a\otimes b\right)   &  =i_{C}\left(  a\right)  \bullet j_{C}\left(
b\right)  =i_{C}\left(  x_{1}\cdots x_{\alpha}\right)  \bullet j_{C}\left(
y_{1}\cdots y_{\beta}\right)  =i_{C}\left(  x_{1}\right)  \bullet\cdots\bullet
i_{C}\left(  x_{\alpha}\right)  \bullet i_{C}\left(  y_{1}\right)
\bullet\cdots\bullet i_{C}\left(  y_{b}\right) \\
&  =\left(  -1\right)  ^{\alpha\beta}i_{C}\left(  y_{1}\right)  \bullet
\cdots\bullet i_{C}\left(  y_{b}\right)  \bullet i_{C}\left(  x_{1}\right)
\bullet\cdots\bullet i_{C}\left(  x_{\alpha}\right)  =\left(  -1\right)
^{\alpha\beta}i_{C}\left(  b\right)  \bullet j_{C}\left(  a\right)  =\left(
-1\right)  ^{\alpha\beta}f\left(  b\otimes a\right)
\end{align}

\end{proof}

One can further show that $f$ is injective on $f^{-1}\left(  \mathcal{C}%
\left(  H\right)  \right)  $ thus an algebra isomorphism from $f^{-1}\left(
\mathcal{C}\left(  H\right)  \right)  \subset\mathcal{C}\left(  E\right)
\otimes\mathcal{C}\left(  F\right)  $
\begin{equation}
a\otimes b=\left(  -1\right)  ^{\alpha\beta}b\otimes a
\end{equation}
so $\otimes$ is a skew tensor product when $a\in f^{-1}\left(  \mathcal{C}%
\left(  E\right)  \right)  $ and $b\in f^{-1}\left(  \mathcal{C}\left(
F\right)  \right)  $ and $\otimes$ is restricted to $f^{-1}\left(
\mathcal{C}\left(  H\right)  \right)  .$

\paragraph{Example}

Let $\left(  \left.  {}\right\vert \right)  $ be degenerate and $E_{0}$ be its
null space and $E_{0}^{c}$ a complimentary subspace so that
\begin{equation}
E=E_{0}\oplus E_{0}^{c}%
\end{equation}
then
\begin{equation}
C\left(  E\right)  \cong\wedge E_{0}\widehat{\otimes}C\left(  E_{0}%
^{c}\right)  =\mathcal{C}\left(  E_{0}\oplus E_{0}^{c}\right)  .
\end{equation}


\subsubsection{\emph{Linear} Isomorphism between Exterior and Clifford
Algebras}

The \emph{linear isomorphism} $\xi_{E}$%
\begin{equation}
\xi_{E}:\bigwedge\left(  E\right)  \rightarrow C\left(  E\right)
\end{equation}
given by
\begin{equation}
\xi_{E}\left(  x_{1}\wedge\cdots\wedge x_{p}\right)  =\frac{1}{p!}\sum
_{\sigma\in S_{p}}\left(  -1\right)  ^{\pi\left(  \sigma\right)  }%
x_{\sigma\left(  1\right)  }\cdots x_{\sigma\left(  p\right)  },
\end{equation}
shows that $\bigwedge\left(  E\right)  $ is linearly identical to $C\left(
E\right)  ,$ though not algebraically identical as%
\begin{equation}
\xi_{E}\left(  A\wedge B\right)  \neq\xi_{E}\left(  A\right)  \circ\xi
_{E}\left(  B\right)  .
\end{equation}
The inverse $\xi_{E}^{-1}$ is given by
\begin{equation}
x_{1}\wedge\cdots\wedge x_{p}=\xi_{E}^{-1}\left(  \frac{1}{p!}\sum_{\sigma\in
S_{p}}\left(  -1\right)  ^{\pi\left(  \sigma\right)  }x_{\sigma\left(
1\right)  }\cdots x_{\sigma\left(  p\right)  }\right)  =\xi_{E}^{-1}\left(
x_{1}\cdots x_{p}\right)  ,
\end{equation}
where $\left\{  \left.  x_{j}\right\vert 1\leq j\leq p\right\}  $ are all
distinct. Again
\begin{equation}
\xi_{E}^{-1}\left(  A\circ B\right)  \neq\xi_{E}^{-1}\left(  A\right)
\wedge\xi_{E}^{-1}\left(  B\right)  .
\end{equation}


Elements of $C\left(  E\right)  $ can be expanded as
\begin{align}
A  &  =\sum_{0\leq p\leq r}\sum_{1\leq j_{1}<\cdots<j_{p}\leq r}\alpha
_{j_{1}\cdots j_{p}}x_{1}\wedge\cdots\wedge x_{p}\\
B  &  =\sum_{0\leq p\leq r}\sum_{1\leq j_{1}<\cdots<j_{p}\leq r}\beta
_{j_{1}\cdots j_{p}}x_{1}\wedge\cdots\wedge x_{p}%
\end{align}


\subsubsection{Centre and Anticentre of $C\left(  E\right)  .$}

The centre, which is \emph{a subalgebra,} is given by
\begin{equation}
Z_{E}=\left\{  A\left\vert \left[  A,B\right]  =0;\forall B\in C\left(
E\right)  ,A\in C\left(  E\right)  \right.  \right\}
\end{equation}
and the anticentre, which is \emph{a subspace,} by%
\begin{equation}
AZ_{E}=\left\{  A\left\vert \left\{  A,B\right\}  =0;\forall B\in C\left(
E\right)  ,A\in C\left(  E\right)  \right.  \right\}  .
\end{equation}
Both the centre and anticentre are $\mathbb{Z}_{2}$-graded and decompose into
even and odd elements
\begin{align}
Z_{E}  &  =Z_{E}^{0}\oplus Z_{E}^{1}\\
AZ_{E}  &  =\left(  AZ_{E}\right)  ^{0}\oplus\left(  AZ_{E}\right)  ^{1}.
\end{align}
Greub shows that if the inner product is non degenerate
\begin{align}
\left(  AZ_{E}\right)  ^{1}  &  =0\\
Z_{E}^{0}  &  =\left\{  I_{A}\right\}  .
\end{align}
If $n$ is even
\begin{align}
Z_{E}  &  =\left\{  I_{A}\right\}  +\left\{  \Delta_{A}\right\} \\
AZ_{E}  &  =0.
\end{align}
If $n$ is odd
\begin{align}
Z_{E}  &  =\left\{  I_{A}\right\} \\
AZ_{E}  &  =\left\{  \Delta_{A}\right\}  .
\end{align}


\begin{theorem}
Let $E$ be an even dimensional space which admits a deterninant function
$\Delta$ such that $e_{\Delta}^{2}=\varepsilon e$ for $\varepsilon=\pm1.$ Then
for any inner product space $F$%
\begin{equation}
C\left(  E\oplus\varepsilon F\right)  \cong C\left(  E\right)  \otimes
C\left(  F\right)
\end{equation}

\end{theorem}

\subsubsection{Direct Sum of Dual Spaces}

Let $E_{1\text{ }}$and $E_{2}$ be dual spaces one can consider the \emph{non
degenerate} inner product
\begin{equation}
\left(  x_{1}\oplus x_{2},y_{1}\oplus y_{2}\right)  =\frac{1}{2}\left(
\left\langle x_{1},y_{2}\right\rangle +\left\langle x_{2},y_{1}\right\rangle
\right)  ;x_{1},y_{1}\in E_{1\text{ }},x_{2},y_{2}\in E_{2\text{ }},
\end{equation}
which is a different inner product to the conventional one i.e.
\begin{equation}
\left(  x_{1}\oplus x_{2},y_{1}\oplus y_{2}\right)  _{con}=\left\langle
x_{1},y_{1}\right\rangle +\left\langle x_{2},y_{2}\right\rangle ;x_{1}%
,y_{1}\in E_{1\text{ }},x_{2},y_{2}\in E_{2\text{ }}.
\end{equation}


\begin{proposition}
$C\left(  E\right)  \cong\mathcal{B}\left(  \wedge\left(  E_{1}\right)
\right)  $
\end{proposition}

\begin{proposition}
Let $E$ be a $2r$-dimensional vector space with a \emph{nondegenerate} inner
product. Assume that there exists an involution $\omega$ of E such that
$\omega^{t}=-\omega$, where the transpose is wrt $\left(  \left.
{}\right\vert \right)  $. Then the Clifford algebra $C\left(  E\right)  $ is
\emph{algebraically} isomorphic to the algebra of linear transformations of
$\wedge\left(  E_{1}\right)  ,$ where
\begin{equation}
E_{1}=\ker\left\{  (\omega-I)\right\}  =\left\{  \left.  \omega\left(
x\right)  =x\right\vert x\in E\right\}  .
\end{equation}

\end{proposition}

\subsubsection{Creation and Annihilation Operators}

Recall the left multiplication operator $\mu\left(  x\right)  $, by the
element $x$, which is proportional to the creation operator $a^{\dagger
}\left(  x\right)  ,$
\begin{align}
\mu\left(  x\right)   &  :\wedge\left(  E_{1}\right)  \rightarrow\wedge\left(
E_{1}\right) \\
\mu\left(  x\right)   &  \in\mathcal{B}\left(  \wedge\left(  E_{1}\right)
\right)
\end{align}
is defined by
\begin{equation}
\left[  \mu\left(  x\right)  \right]  \left(  y\right)  =x\wedge y.
\end{equation}
Since $\wedge\left(  E_{1}\right)  $ is an associative algebra the
relationship between the exterior product and the operator/composition product
is
\begin{equation}
\mu\left(  x\wedge y\right)  =\mu\left(  x\right)  \circ\mu\left(  y\right)
\end{equation}
and
\begin{equation}
\mu\left(  x\right)  \circ\mu\left(  y\right)  =\left(  -1\right)  ^{pq}%
\mu\left(  y\right)  \circ\mu\left(  x\right)  .
\end{equation}
The adjoint operator (also associated with $x\in\wedge\left(  E_{1}\right)
)$
\begin{align}
i\left(  x\right)   &  :\wedge\left(  E_{1}\right)  ^{d}\leftarrow
\wedge\left(  E_{1}\right)  ^{d}\\
\mu\left(  x\right)   &  \in\mathcal{B}\left(  \wedge\left(  E_{1}\right)
^{d}\right)
\end{align}
is given by%
\begin{equation}
\left\langle \left.  i\left(  x\right)  y^{d}\right\vert z\right\rangle
=\left\langle \left.  y^{d}\right\vert \mu\left(  x\right)  z\right\rangle
\end{equation}
is proportional to the annihilation operator $a\left(  x\right)  .$

In this case the relationship between the exterior product and the
operator/composition product is%
\begin{equation}
i\left(  x\wedge y\right)  =i\left(  y\right)  \circ i\left(  x\right)
\end{equation}
and
\begin{equation}
i\left(  x\right)  \circ i\left(  y\right)  =\left(  -1\right)  ^{pq}i\left(
y\right)  \circ i\left(  x\right)  ,
\end{equation}
where $x,y\in\wedge\left(  E_{1}\right)  .$

Let $x,y\in E_{1}$ then the operators $i\left(  x\right)  ,i\left(  y\right)
$ are homogeneous of degree $-1$ and
\begin{equation}
i\left(  x\right)  \circ i\left(  y\right)  +i\left(  y\right)  \circ i\left(
x\right)  =0.
\end{equation}
In particular they are nilpotent i.e
\begin{equation}
i\left(  x\right)  ^{2}=0.
\end{equation}


\begin{proposition}
The operator $i\left(  x\right)  $ is an anti derivation in the algebra
$\wedge\left(  E_{1}\right)  ^{d}\equiv\wedge\left(  E_{1}^{d}\right)  $ i.e.
\begin{equation}
\left[  i\left(  x\right)  \right]  \left(  u^{d}\wedge v^{d}\right)  =\left(
\left[  i\left(  x\right)  \right]  \left(  u^{d}\right)  \right)  \wedge
v^{d}+\left(  -1\right)  ^{p}u^{d}\wedge\left(  \left[  i\left(  x\right)
\right]  \left(  v^{d}\right)  \right)  ;u^{d}\in\wedge^{p}\left(  E_{1}%
^{d}\right)  ,v^{d}\in\wedge\left(  E_{1}^{d}\right)  .
\end{equation}

\end{proposition}

\begin{corollary}
1: CAR%
\begin{equation}
i\left(  x\right)  \circ\mu\left(  y^{d}\right)  +\mu\left(  y^{d}\right)
\circ i\left(  x\right)  =\left\langle \left.  y^{d}\right\vert x\right\rangle
I_{\mathcal{A}}%
\end{equation}

\end{corollary}

\begin{corollary}
II:%
\begin{equation}
\left[  i\left(  x\right)  \right]  \left(  y_{1}^{d}\wedge\cdots\wedge
y_{p}^{d}\right)  =\sum_{1\leq j\leq p}\left(  -1\right)  ^{p-1}\left\langle
\left.  y_{j}^{d}\right\vert x\right\rangle y_{1}^{d}\wedge\cdots
\wedge\widehat{y_{j}^{d}}\wedge\cdots\wedge y_{p}^{d}%
\end{equation}

\end{corollary}

\begin{corollary}
III:Let $F_{1}$ be subspace of $E_{1}$ and $F_{1}^{\bot}$ its orthogonal
compliment then $\wedge\left(  F_{1}\right)  $ and $\wedge\left(  F_{1}%
^{d}\right)  $ are subalgebras of $\wedge\left(  E_{1}\right)  $ and
$\wedge\left(  E_{1}^{d}\right)  $ respectively then
\begin{equation}
\left[  i\left(  x\right)  \right]  \left(  u^{d}\wedge v^{d}\right)  =\left(
-1\right)  ^{pq}u^{d}\wedge\left(  \left[  i\left(  x\right)  \right]  \left(
v^{d}\right)  \right)  ;\,x\in\wedge^{p}\left(  F_{1}\right)  ,u^{d}\in
\wedge^{q}\left(  F_{1}^{\bot}\right)  ,v^{d}\in\wedge\left(  E_{1}%
^{d}\right)  .
\end{equation}

\end{corollary}

\begin{proposition}
If $u^{d}\in\wedge^{p}\left(  E_{1}^{d}\right)  $ satisfies the equation
\begin{equation}
\left[  i\left(  x\right)  \right]  u^{d}=0
\end{equation}
for all $x\in E_{1}$ then
\begin{equation}
u^{d}=0.
\end{equation}

\end{proposition}

In the following $1\leq j<k\leq r$
\begin{align}
x  &  \in E_{1};a\in\wedge^{p}\left(  E_{1}\right) \\
x_{j}^{2}  &  =I_{\mathcal{A}}.
\end{align}


\begin{lemma}
$\xi_{E}\left(  a\wedge x\right)  =\xi_{E}\left(  a\right)  \circ x-\xi
_{E}\left(  \left[  i\left(  x\right)  \right]  \left(  a\right)  \right)  $
\end{lemma}

\begin{lemma}
$\xi_{E}\left(  x\wedge a\right)  =x\circ\xi_{E}\left(  a\right)  -\left(
-1\right)  ^{p}\xi_{E}\left(  \left[  i\left(  x\right)  \right]  \left(
a\right)  \right)  $
\end{lemma}

Examples:

\begin{enumerate}
\item
\begin{align}
a  &  =x_{j}\\
x  &  =x_{k}%
\end{align}%
\begin{align}
\xi_{E}\left(  x_{j}\wedge x_{k}\right)   &  =\xi_{E}\left(  x_{j}\right)
\circ x_{k}-\xi_{E}\left(  i\left(  x_{k}\right)  \circ x_{j}\right) \\
&  =x_{j}x_{k}-\delta_{jk}x_{j}^{2}=\left\{
\begin{array}
[c]{c}%
x_{j}x_{k};\,j\neq k\\
0;\,j=k
\end{array}
\right.
\end{align}


\item
\begin{align}
a  &  =x_{j}x_{k}\\
x  &  =x_{l}%
\end{align}%
\begin{align}
&  \xi_{E}\left(  x_{j}x_{k}\wedge x_{l}\right) \\
&  =\xi_{E}\left(  x_{j}x_{k}\right)  x_{l}-\xi_{E}\left(  \left[  i\left(
x_{l}\right)  \right]  \left(  x_{j}x_{k}\right)  \right) \\
&  =\left(  x_{j}\wedge x_{k}\right)  x_{l}-\xi_{E}\left(  \left[  i\left(
x_{l}\right)  \right]  \left(  x_{j}x_{k}\right)  \right) \\
&  =\frac{1}{2}\left(  x_{j}\wedge x_{k}x_{l}-x_{k}\wedge x_{j}x_{l}\right)
-\xi_{E}\left(  \left(  \left(  \left[  i\left(  x_{l}\right)  \right]
\left(  x_{j}\right)  \right)  x_{k}-x_{j}\left[  i\left(  x_{l}\right)
\right]  \left(  x_{k}\right)  \right)  \right) \\
&  =x_{j}x_{k}x_{l}+x_{j}\delta_{kl}-x_{k}\delta_{jl}%
\end{align}

\end{enumerate}

In the following $u\in C\left(  E\right)  ,x\in E_{1}$

\begin{lemma}
$\left[  \xi_{E}^{-1}\right]  \left(  xu\right)  =x\wedge\left[  \xi_{E}%
^{-1}\right]  \left(  u\right)  -\left[  i\left(  x\right)  \right]  \left(
\xi_{E}^{-1}\left(  u\right)  \right)  $
\end{lemma}

\begin{lemma}
$\left[  \xi_{E}^{-1}\right]  \left(  ux\right)  =\left[  \xi_{E}^{-1}\left(
u\right)  \right]  \wedge x-\left[  i\left(  x\right)  \right]  \left(
\xi_{E}^{-1}\left(  \omega_{E_{1}}\left(  u\right)  \right)  \right)  $
\end{lemma}

\section{Complex Clifford Algebras}

In this case $\left(  \left.  {}\right\vert \right)  $ is a symmetric complex
valued non degenerate bilinear form on $\mathcal{H}^{1}\cong\mathbb{C}^{r}$

\subparagraph{Example: $r=1,$ $C_{1}\protect\cong\mathbb{C\oplus C}$}

Here the isomorphism is given by
\begin{align}
\Phi &  :C_{1}\rightarrow\mathbb{C\oplus C}\\
\Phi\left(  \alpha e_{1}+\beta e_{2}\right)   &  =\left(  \alpha+\beta
,\alpha-\beta\right)  .
\end{align}


A normed determinant function in $\mathbb{C}^{r}$ is given by
\begin{equation}
\Delta\left(  x_{1},\cdots,x_{r}\right)  \Delta\left(  y_{1},\cdots
,y_{r}\right)  =\det\left\{  \left(  x_{i},y_{j}\right)  \right\}
;x_{i},y_{j}\in\mathbb{C}^{r}%
\end{equation}
The cannonical element $e_{\Delta}$ is given by
\begin{equation}
e_{\Delta}^{2}=\left(  -1\right)  ^{r\left(  r-1\right)  }e,
\end{equation}
and one has

\begin{theorem}
$C_{2m+k}\cong C_{2m}\otimes C_{k}.$
\end{theorem}

\begin{proposition}
$C_{2m}\cong\mathcal{B}\left(  \wedge\mathbb{C}^{m}\right)  \cong%
\mathcal{B}\left(  \wedge E_{1}\right)  .$
\end{proposition}

\begin{proposition}
$C_{2m+1}\cong\mathcal{B}\left(  \wedge\mathbb{C}^{m}\right)  \oplus
\mathcal{B}\left(  \wedge\mathbb{C}^{m}\right)  \cong\mathcal{B}\left(  \wedge
E_{1}\right)  \oplus\mathcal{B}\left(  \wedge E_{1}\right)  .$
\end{proposition}

\begin{proof}
Let $\left\{  \left.  a_{j}\right\vert 1\leq j\leq m\right\}  $ and $\left\{
\left.  b_{j}\right\vert 1\leq j\leq m\right\}  $ be mutually orthogonal conb
bases of $A$ and $B$ respectively so that
\begin{align}
A &  =\mathbb{C-}\operatorname*{linspan}\left\{  \left.  a_{j}\right\vert
1\leq j\leq m\right\}  \\
B &  =\mathbb{C-}\operatorname*{linspan}\left\{  \left.  b_{j}\right\vert
1\leq j\leq m\right\}  =A^{\bot},
\end{align}
so that
\begin{equation}
E=A\oplus B
\end{equation}
and
\begin{align}
\left(  \left.  a_{j}\right\vert b_{k}\right)   &  =0\\
\left(  \left.  a_{j}\right\vert a_{k}\right)   &  =\delta_{jk}\\
\left(  \left.  b_{j}\right\vert b_{k}\right)   &  =\delta_{jk}.
\end{align}
\emph{Define an antisymmetric involution} $\omega$
\begin{equation}
\omega:E\rightarrow E
\end{equation}
by%
\begin{align}
\omega\left(  a_{j}\right)   &  =ib_{j}\\
\omega\left(  b_{j}\right)   &  =-ia_{j}.
\end{align}
This is an \emph{involution} as
\begin{align}
\omega\left(  \omega\left(  a_{j}\right)  \right)   &  =i\omega\left(
b_{j}\right)  =-i^{2}a_{j}=a_{j}\\
\omega\left(  \omega\left(  b_{j}\right)  \right)   &  =-i\omega\left(
a_{j}\right)  =-i^{2}b_{j}=b_{j}%
\end{align}
and it is \emph{antisymmetric} as
\begin{align}
\left(  \left.  a_{j}\right\vert \omega\left(  a_{k}\right)  \right)   &
=\left(  \left.  b_{j}\right\vert \omega\left(  b_{k}\right)  \right)  =0\\
\left(  \left.  a_{j}\right\vert \omega\left(  b_{k}\right)  \right)   &
=-i\left(  \left.  a_{j}\right\vert a_{k}\right)  =-i\delta_{jk}\\
\left(  \left.  b_{j}\right\vert \omega\left(  a_{k}\right)  \right)   &
=i\left(  \left.  b_{j}\right\vert b_{k}\right)  =i\delta_{jk}.
\end{align}
The involution $\omega$ in turn defines a subspace of $E_{1}$ of $E$ by
\begin{equation}
E_{1}=\left\{  \left.  x\right\vert \omega\left(  x\right)  =x;x\in E\right\}
=\ker\left\{  \omega-I_{\mathcal{A}}\right\}  ,
\end{equation}
so that
\begin{align}
E_{1} &  =\mathbb{C-}\operatorname*{linspan}\left\{  \left.  e_{j}\right\vert
1\leq j\leq m\right\}  \\
E_{1}^{\bot} &  =\mathbb{C-}\operatorname*{linspan}\left\{  \left.
e_{j}\right\vert m+1\leq j\leq2m\right\}  ,
\end{align}
where
\begin{align}
e_{j} &  =\frac{1}{\sqrt{2}}\left(  a_{j}+ib_{j}\right)  ;\,1\leq j\leq m\\
e_{j} &  =\frac{1}{\sqrt{2}}\left(  a_{j}-ib_{j}\right)  ;m+\,1\leq j\leq2m.
\end{align}
The effect of $\omega$ on the basis $\left\{  \left.  e_{j}\right\vert 1\leq
j\leq2m\right\}  ,\left\{  \left.  a_{j}\right\vert 1\leq j\leq m\right\}  $
and $\left\{  \left.  b_{j}\right\vert 1\leq j\leq m\right\}  $ is given by
\begin{align}
\omega\left(  e_{j}\right)   &  =\frac{1}{\sqrt{2}}\omega\left(  a_{j}%
+ib_{j}\right)  =\frac{1}{\sqrt{2}}\left(  ib_{j}-iia_{j}\right)  =\,\frac
{1}{\sqrt{2}}\left(  a_{j}+ib_{j}\right)  \,\,=e_{j};\,1\leq j\leq m\\
\omega\left(  e_{j}\right)   &  =\frac{1}{\sqrt{2}}\omega\left(  a_{j}%
-ib_{j}\right)  =\frac{1}{\sqrt{2}}\left(  ib_{j}+iia_{j}\right)  =-\frac
{1}{\sqrt{2}}\left(  a_{j}-ib_{j}\right)  =-e_{j};\,m+\,1\leq j\leq2m.
\end{align}
The representations of $A$ and $B$ wrt the bases $\left\{  \left.
a_{j}\right\vert 1\leq j\leq m\right\}  $,$\left\{  \left.  b_{j}\right\vert
1\leq j\leq m\right\}  ,\left\{  \left.  e_{j}\right\vert 1\leq j\leq
m\right\}  $ and $\left\{  \left.  e_{j}\right\vert m+1\leq j\leq2m\right\}  $
give
\begin{equation}
\mathbb{C}^{2m}=\mathbb{A\oplus B=E\oplus E}^{\bot}=\mathbb{E\oplus}%
\overline{\mathbb{E}}?,
\end{equation}
where
\begin{align}
\mathbb{A} &  =\mathcal{R}_{a}\left(  A\right)  \\
\mathbb{B} &  =\mathcal{R}_{b}\left(  B\right)  ,
\end{align}%
\begin{align}
\mathbb{E} &  =\mathcal{R}e\left(  E\right)  \\
\mathbb{E}^{\bot} &  =\overline{\mathbb{E}}=\mathcal{R}_{e}\left(  E^{\bot
}\right)  ,?
\end{align}
and the complex conjugation is defined wrt the basis $\left\{  \left.
e_{j}\right\vert 1\leq j\leq2m\right\}  .$

The representations of $\omega$ are given by
\begin{equation}
\omega\left(  \mathbf{a,b}\right)  =\left(  \mathbf{a,b}\right)
\mathcal{R}_{ab}\left(  \omega\right)  =\left(  \mathbf{a,b}\right)  \left(
\begin{array}
[c]{cc}%
\mathbb{O} & -i\mathbb{I}_{m}\\
i\mathbb{I}_{m} & \mathbb{O}%
\end{array}
\right)  =\left(  \mathbf{a,b}\right)  i\left(
\begin{array}
[c]{cc}%
\mathbb{O} & -\mathbb{I}_{m}\\
\mathbb{I}_{m} & \mathbb{O}%
\end{array}
\right)
\end{equation}
and%
\begin{equation}
\omega\left(  \mathbf{e}_{1}\mathbf{,e}_{2}\right)  =\left(  \mathbf{e}%
_{1}\mathbf{,e}_{2}\right)  \mathcal{R}_{e}\left(  \omega\right)  =\left(
\mathbf{e}_{1}\mathbf{,e}_{2}\right)  \left(
\begin{array}
[c]{cc}%
\mathbb{I}_{m} & \mathbb{O}\\
\mathbb{O} & -\mathbb{I}_{m}%
\end{array}
\right)  .
\end{equation}%
\begin{equation}
\left(  \mathbf{a,b}\right)  \frac{1}{\sqrt{2}}\left(
\begin{array}
[c]{cc}%
\mathbb{I}_{m} & \mathbb{I}_{m}\\
i\mathbb{I}_{m} & -i\mathbb{I}_{m}%
\end{array}
\right)  =\left(  \mathbf{e}_{1}\mathbf{,e}_{2}\right)
\end{equation}

\end{proof}

$\lceil$Example $m=1$%
\begin{equation}
\mathbb{C}^{2}=\mathbb{A\oplus B}%
\end{equation}%
\begin{align}
\mathbb{A}=\mathbb{C-}  &  \operatorname*{linspan}\left\{  a_{1}\right\} \\
\mathbb{B}=\mathbb{C-}  &  \operatorname*{linspan}\left\{  b_{1}\right\}
\end{align}
That
\begin{align}
E_{1}  &  =\mathbb{C-}\operatorname*{linspan}\left\{  a_{1}+ib_{1}\right\} \\
E_{1}^{c}  &  =\mathbb{C-}\operatorname*{linspan}\left\{  a_{1}-ib_{1}%
\right\}
\end{align}
is shown by
\begin{align}
\omega\left(  a_{1}\right)   &  =ib_{1}\\
\omega\left(  b_{1}\right)   &  =-ia_{1}%
\end{align}
giving%
\begin{equation}
\omega\left(  \alpha a_{1}+\beta b_{1}\right)  =i\alpha b_{1}-i\beta
a_{1};\alpha,\beta\in\mathbb{C}.
\end{equation}
The eigenvlaue $+1$ constraint gives
\[
i\alpha b_{1}-i\beta a_{1}=\alpha a_{1}+\beta b_{1}%
\]
leading to
\begin{align}
i\alpha &  =\beta\\
-i\beta &  =\alpha\Leftrightarrow i\alpha=\beta
\end{align}
and thus
\begin{equation}
\alpha a_{1}+\beta b_{1}=\alpha a_{1}+i\alpha b_{1}=\alpha\left(  a_{1}%
+ib_{1}\right)  ,
\end{equation}
so the $+1$ eigenspace is populated by vectors of the form $\left\{  \left.
\alpha\left(  a_{1}+ib_{1}\right)  \right\vert \alpha\in\mathbb{C}\right\}  .$

The eigenvlaue $-1$ constraint gives
\[
i\alpha b_{1}-i\beta a_{1}=-\alpha a_{1}-\beta b_{1}%
\]
leading to
\begin{align}
i\alpha &  =-\beta\\
i\beta &  =\alpha\Leftrightarrow i\alpha=-\beta
\end{align}
and thus
\begin{equation}
\alpha a_{1}+\beta b_{1}=\alpha a_{1}-i\alpha b_{1}=\alpha\left(  a_{1}%
-ib_{1}\right)  ,
\end{equation}
so the $-1$ eigenspace is populated by vectors of the form $\left\{  \left.
\alpha\left(  a_{1}-ib_{1}\right)  \right\vert \alpha\in\mathbb{C}\right\}  .$

$\lceil$%
\begin{align}
\left(  x+iy\right)  \left(  a_{1}+ib_{1}\right)   &  =xa_{1}+iya_{1}%
+xib_{1}-yb_{1}=\left(  x+iy\right)  a_{1}+\left(  ix-y\right)  b_{1}=\left(
x+iy\right)  a_{1}+i\left(  x+iy\right)  b_{1}=\left(  x+iy\right)  \left(
a_{1}+ib_{1}\right) \\
\left(  x+iy\right)  \left(  a_{1}-ib_{1}\right)   &  =xa_{1}+iya_{1}%
-xib_{1}+yb_{1}=\left(  x+iy\right)  a_{1}+\left(  -ix+y\right)  b_{1}=\left(
x+iy\right)  a_{1}-i\left(  x+iy\right)  b_{1}=\left(  x+iy\right)  \left(
a_{1}-ib_{1}\right)  .
\end{align}
$\rfloor$

$A$ could be the \emph{complex} spaces generated by $\left\{  \left.
a_{j}^{\dagger}\right\vert 1\leq j\leq m\right\}  $ or $\left\{  \left.
x_{j}\right\vert 1\leq j\leq m\right\}  $ and $B$ by $\left\{  \left.
a_{j}\right\vert 1\leq j\leq m\right\}  $ or $\left\{  \left.  x_{j}%
\right\vert m+1\leq j\leq2m\right\}  $ then%

\begin{align}
&  \left(  \operatorname{Re}\mathbf{z,}\operatorname{Im}\mathbf{z}\right)
\frac{1}{\sqrt{2}}\left(
\begin{array}
[c]{cc}%
\mathbb{I}_{m} & \mathbb{I}_{m}\\
i\mathbb{I}_{m} & -i\mathbb{I}_{m}%
\end{array}
\right)  =\left(  \mathbf{z,}\overset{-}{\mathbf{z}}\right)  \\
&  \left(  \operatorname{Re}\mathbf{z,}\operatorname{Im}\mathbf{z}\right)
=\left(  \mathbf{z,}\overset{-}{\mathbf{z}}\right)  \frac{1}{\sqrt{2}}\left(
\begin{array}
[c]{cc}%
\mathbb{I}_{m} & -i\mathbb{I}_{m}\\
\mathbb{I}_{m} & i\mathbb{I}_{m}%
\end{array}
\right)
\end{align}
and%
\begin{equation}
\left(  \operatorname{Re}\mathbf{z,}\operatorname{Im}\mathbf{z}\right)
\frac{1}{\sqrt{2}}\left(
\begin{array}
[c]{cc}%
\mathbb{I}_{m} & \mathbb{I}_{m}\\
i\mathbb{I}_{m} & -i\mathbb{I}_{m}%
\end{array}
\right)  \left(
\begin{array}
[c]{cc}%
\mathbb{O}_{m} & \mathbb{I}_{m}\\
\mathbb{I}_{m} & \mathbb{O}_{m}%
\end{array}
\right)  =\left(  \operatorname{Re}\mathbf{z,}\operatorname{Im}\mathbf{z}%
\right)  \frac{1}{\sqrt{2}}\left(
\begin{array}
[c]{cc}%
\mathbb{I}_{m} & \mathbb{I}_{m}\\
-i\mathbb{I}_{m} & i\mathbb{I}_{m}%
\end{array}
\right)  =\left(  \mathbf{z,}\overset{-}{\mathbf{z}}\right)  \left(
\begin{array}
[c]{cc}%
\mathbb{O}_{m} & \mathbb{I}_{m}\\
\mathbb{I}_{m} & \mathbb{O}_{m}%
\end{array}
\right)  =\left(  \overset{-}{\mathbf{z}},\mathbf{z}\right)  .
\end{equation}
Complex conjugation is given by
\begin{align}
&  \left(  \operatorname{Re}\mathbf{z,}\operatorname{Im}\mathbf{z}\right)
\left(
\begin{array}
[c]{cc}%
\mathbb{I}_{m} & \mathbb{O}_{m}\\
\mathbb{O}_{m} & -\mathbb{I}_{m}%
\end{array}
\right)  =\left(  \operatorname{Re}\mathbf{z,-}\operatorname{Im}%
\mathbf{z}\right)  \\
&  =\left(  \operatorname{Re}\mathbf{z,}\operatorname{Im}\mathbf{z}\right)
\frac{1}{\sqrt{2}}\left(
\begin{array}
[c]{cc}%
\mathbb{I}_{m} & \mathbb{I}_{m}\\
i\mathbb{I}_{m} & -i\mathbb{I}_{m}%
\end{array}
\right)  \frac{1}{\sqrt{2}}\left(
\begin{array}
[c]{cc}%
\mathbb{I}_{m} & -i\mathbb{I}_{m}\\
\mathbb{I}_{m} & i\mathbb{I}_{m}%
\end{array}
\right)  \left(
\begin{array}
[c]{cc}%
\mathbb{I}_{m} & \mathbb{O}_{m}\\
\mathbb{O}_{m} & -\mathbb{I}_{m}%
\end{array}
\right)  \frac{1}{\sqrt{2}}\left(
\begin{array}
[c]{cc}%
\mathbb{I}_{m} & \mathbb{I}_{m}\\
i\mathbb{I}_{m} & -i\mathbb{I}_{m}%
\end{array}
\right)  \\
&  =\left(  \operatorname{Re}\mathbf{z,}\operatorname{Im}\mathbf{z}\right)
\frac{1}{\sqrt{2}}\left(
\begin{array}
[c]{cc}%
\mathbb{I}_{m} & \mathbb{I}_{m}\\
i\mathbb{I}_{m} & -i\mathbb{I}_{m}%
\end{array}
\right)  \frac{1}{2}\left(
\begin{array}
[c]{cc}%
\mathbb{I}_{m} & i\mathbb{I}_{m}\\
\mathbb{I}_{m} & -i\mathbb{I}_{m}%
\end{array}
\right)  \left(
\begin{array}
[c]{cc}%
\mathbb{I}_{m} & \mathbb{I}_{m}\\
i\mathbb{I}_{m} & -i\mathbb{I}_{m}%
\end{array}
\right)  \\
&  =\left(  \operatorname{Re}\mathbf{z,}\operatorname{Im}\mathbf{z}\right)
\frac{1}{\sqrt{2}}\left(
\begin{array}
[c]{cc}%
\mathbb{I}_{m} & \mathbb{I}_{m}\\
i\mathbb{I}_{m} & -i\mathbb{I}_{m}%
\end{array}
\right)  \left(
\begin{array}
[c]{cc}%
\mathbb{O}_{m} & \mathbb{I}_{m}\\
\mathbb{I}_{m} & \mathbb{O}_{m}%
\end{array}
\right)  \\
&  =\left(  \overset{-}{\mathbf{z}},\mathbf{z}\right)  \left(
\begin{array}
[c]{cc}%
\mathbb{O}_{m} & \mathbb{I}_{m}\\
\mathbb{I}_{m} & \mathbb{O}_{m}%
\end{array}
\right)  =\left(  \overset{-}{\mathbf{z}},\mathbf{z}\right)
\end{align}


\subsection{Complexification of Real Clifford Algebra}

Let $F$ be a real inner product space (not neccesarily finite dimension) its
complexification
\begin{equation}
E=\mathbb{C\otimes}F,
\end{equation}
with inner product
\begin{equation}
\left(  \left.  \lambda\otimes x\right\vert \mu\otimes y\right)  =\lambda
\mu\left(  \left.  x\right\vert y\right)
\end{equation}
then one can show
\begin{equation}
C\left(  E\right)  \cong\mathbb{C\otimes}C\left(  F\right)  .
\end{equation}


\subsection{Clifford Algebras Over a Real Vector Space}

If $\left(  \left.  x\right\vert y\right)  $ is a {\footnotesize non
degenerate} inner prodct on $E$ then
\begin{equation}
E=E_{+}\oplus E_{-},
\end{equation}
where
\begin{align}
\dim\left\{  E\right\}   &  =n=p+q\\
\dim\left\{  E_{+}\right\}   &  =p\\
\dim\left\{  E_{-}\right\}   &  =q.
\end{align}


We write
\begin{align}
C\left(  p,0\right)   &  =C\left(  E_{+}\right) \\
C\left(  0,q\right)   &  =C\left(  E_{-}\right)
\end{align}


\subsubsection{Canonical Element}

The canonical element of $C\left(  p,q\right)  $ is given by
\begin{equation}
e_{\Delta}^{2}=\left(  -1\right)  ^{\frac{1}{2}n\left(  n-1\right)  +q}e,
\end{equation}
where
\begin{equation}
e=e_{1}\wedge\cdots\wedge e_{n}.
\end{equation}


\begin{enumerate}
\item If $\left(  \left.  x\right\vert y\right)  $ is positive definite
\begin{equation}
e_{\Delta}^{2}=\left(  -1\right)  ^{\frac{1}{2}n\left(  n-1\right)  }e
\end{equation}


\item If $\left(  \left.  x\right\vert y\right)  $ is negative definite
\begin{equation}
e_{\Delta}^{2}=\left(  -1\right)  ^{\frac{1}{2}n\left(  n-1\right)  }e
\end{equation}


\item If $n=2m$ and $p=q=m,$
\begin{equation}
e_{\Delta}^{2}=e.
\end{equation}

\end{enumerate}

\begin{theorem}
There are algebra isomorphisms
\end{theorem}

%

\begin{align}
C\left(  p,q\right)  \mathbb{\otimes}C\left(  2,0\right)   &  \cong C\left(
q+2,p\right) \\
C\left(  p,q\right)  \mathbb{\otimes}C\left(  0,2\right)   &  \cong C\left(
q,p+2\right)
\end{align}
In particular for $n\geq0$
\begin{align}
C\left(  0,n\right)  \mathbb{\otimes}C\left(  2,0\right)   &  \cong C\left(
n+2,0\right) \\
C\left(  n,0\right)  \mathbb{\otimes}C\left(  0,2\right)   &  \cong C\left(
0,n+2\right)
\end{align}


\begin{theorem}
Assume that $s\equiv0\,\left(  \text{mod }4\right)  $ then
\end{theorem}

%

\begin{equation}
C\left(  p,q\right)  \cong C\left(  q,p\right)
\end{equation}
In particular
\begin{equation}
C\left(  0,n\right)  \cong C\left(  n,0\right)
\end{equation}
if $n\equiv0\,\left(  \text{mod }4\right)  .$

\subsection{Inner Product Spaces with Zero Signiture}

If the inner product $\left(  \left.  {}\right\vert \right)  $ is of type
$\left(  m,m\right)  $ it has zero signiture and one has the following theorem:

\begin{theorem}%
\begin{equation}
C\left(  E\right)  \cong\mathcal{B}\left(  \wedge\left(  E_{1}\right)
\right)  ,
\end{equation}
where
\begin{equation}
\dim\left\{  E_{1}\right\}  =m.
\end{equation}

\end{theorem}

\begin{proof}
Let
\begin{align}
E_{+}  &  =\mathbb{C-}\operatorname*{linspan}\left\{  \left.  a_{j}\right\vert
1\leq j\leq m\right\} \\
E_{-}  &  =\mathbb{C-}\operatorname*{linspan}\left\{  \left.  b_{j}\right\vert
1\leq j\leq m\right\}  =E_{+}^{\bot},
\end{align}
so that
\begin{equation}
E=E_{+}\oplus E_{-}%
\end{equation}
where
\begin{align}
\left(  \left.  a_{j}\right\vert b_{k}\right)   &  =0\\
\left(  \left.  a_{j}\right\vert a_{k}\right)   &  =\delta_{jk}\\
\left(  \left.  b_{j}\right\vert b_{k}\right)   &  =-\delta_{jk}.
\end{align}
Define an involution $\omega$ of $E$ by
\begin{align}
\omega\left(  a_{j}\right)   &  =b_{j}\\
\omega\left(  b_{j}\right)   &  =a_{j};1\leq j\leq m.
\end{align}
Then $\omega$ is antisymmetric wrt inner product $\left(  \left.
{}\right\vert \right)  $ and the $m$ dimensional subspace $E_{1}$ defined by
\begin{equation}
E_{1}=\ker\left\{  \omega-I_{E}\right\}  =\mathbb{R-}\operatorname*{linspan}%
\left\{  \left.  x\in E\right\vert \omega\left(  x\right)  =x\right\}
\end{equation}
generates the Clifford algebra $C\left(  E\right)  $ by the vector space
isomorphism
\begin{equation}
C\left(  E\right)  \cong\mathcal{B}\left(  \wedge\left(  E_{1}\right)
\right)  .
\end{equation}
The subspaces $E_{1}$ and $E_{1}^{c}$ are given by elements of the form
\begin{align}
\omega\left(  x+\omega\left(  x\right)  \right)   &  =\omega\left(  x\right)
+x=x+\omega\left(  x\right) \\
\omega\left(  x-\omega\left(  x\right)  \right)   &  =-\omega\left(  x\right)
+x=-\left(  x+\omega\left(  x\right)  \right)
\end{align}
Thus%
\begin{equation}
E_{1}=\mathbb{R-}\operatorname*{linspan}\left\{  \left.  \frac{1}{\sqrt{2}%
}\left(  a_{j}+b_{j}\right)  \right\vert 1\leq j\leq m\right\}
\end{equation}
\begin{equation}
E_{1}^{c}=\mathbb{R-}\operatorname*{linspan}\left\{  \left.  \frac{1}{\sqrt
{2}}\left(  a_{j}-b_{j}\right)  \right\vert 1\leq j\leq m\right\}  ,
\end{equation}
which shows that the subspaces $E_{1}$ and $E_{1}^{c}$ are \emph{totally
isotropic}, but not orthogonal to each other as
\begin{equation}
\left(  \left.  \frac{1}{\sqrt{2}}\left(  a_{j}+b_{j}\right)  \right\vert
\frac{1}{\sqrt{2}}\left(  a_{k}+b_{k}\right)  \right)  =\frac{1}{2}\left\{
\left(  \left.  a_{j}\right\vert a_{k}\right)  +\left(  \left.  a_{j}%
\right\vert b_{k}\right)  +\left(  \left.  b_{j}\right\vert a_{k}\right)
+\left(  \left.  b_{j}\right\vert b_{k}\right)  \right\}  =0;1\leq j,k\leq m
\end{equation}%
\begin{equation}
\left(  \left.  \frac{1}{\sqrt{2}}\left(  a_{j}-b_{j}\right)  \right\vert
\frac{1}{\sqrt{2}}\left(  a_{k}-b_{k}\right)  \right)  =\frac{1}{2}\left\{
\left(  \left.  a_{j}\right\vert a_{k}\right)  -\left(  \left.  a_{j}%
\right\vert b_{k}\right)  -\left(  \left.  b_{j}\right\vert a_{k}\right)
+\left(  \left.  b_{j}\right\vert b_{k}\right)  \right\}  =0;1\leq j,k\leq m
\end{equation}%
\begin{equation}
\left(  \left.  \frac{1}{\sqrt{2}}\left(  a_{j}+b_{k}\right)  \right\vert
\frac{1}{\sqrt{2}}\left(  a_{j}-b_{k}\right)  \right)  =\frac{1}{2}\left\{
\left(  \left.  a_{j}\right\vert a_{k}\right)  -\left(  \left.  a_{j}%
\right\vert b_{k}\right)  +\left(  \left.  b_{j}\right\vert a_{k}\right)
-\left(  \left.  b_{j}\right\vert b_{k}\right)  \right\}  =\delta_{jk};1\leq
j,k\leq m.
\end{equation}
The real subspaces $E_{1}$ and $E_{1}^{c}$ correspond to real linear
combinations of creators and annihilators respectiely, while the real
subspaces $E_{+}$ and $E_{-}$ correspond to real linear combinations of
Majorona operators.
\end{proof}

\subsection{Definite Real Clifford Algebras of Low Dimension}

The clifford algebras $C\left(  n,0\right)  $ and $C\left(  0,n\right)  $ for
$1\leq n\leq8$ are given by
\begin{equation}
C\left(  1,0\right)  \cong\mathbb{R\oplus R}%
\end{equation}%
\begin{equation}
C\left(  0,1\right)  \cong\mathbb{C}%
\end{equation}%
\begin{equation}
C\left(  2,0\right)  \cong\mathcal{B}\left(  \mathbb{R}^{2}\right)
\end{equation}%
\begin{equation}
C\left(  3,0\right)  \cong\mathbb{C}\otimes\mathcal{B}\left(  \mathbb{R}%
^{2}\right)  =M\left(  2,\mathbb{C}\right)
\end{equation}%
\begin{equation}
C\left(  0,3\right)  \cong\mathbb{H\oplus H}%
\end{equation}%
\begin{equation}
C\left(  4,0\right)  \cong C\left(  0,4\right)  \cong\mathbb{H}\otimes
\mathcal{B}\left(  \mathbb{R}^{2}\right)  =M\left(  2,\mathbb{H}\right)
\end{equation}%
\begin{equation}
C\left(  5,0\right)  \cong\left(  \mathbb{H\oplus H}\right)  \otimes
\mathcal{B}\left(  \mathbb{R}^{2}\right)  \cong\mathbb{H}\otimes
\mathcal{B}\left(  \mathbb{R}^{2}\right)  \oplus\mathcal{B}\left(
\mathbb{R}^{2}\right)  \otimes\mathbb{H}=M\left(  2,\mathbb{H}\right)  \oplus
M\left(  2,\mathbb{H}\right)
\end{equation}%
\begin{equation}
C\left(  0,5\right)  \cong\left(  \mathbb{H\oplus H}\right)  \otimes
\mathcal{B}\left(  \mathbb{R}^{2}\right)  \cong\mathbb{C}\otimes
\mathcal{B}\left(  \mathbb{R}^{2}\right)  \otimes\mathbb{H}%
\end{equation}%
\begin{equation}
C\left(  6,0\right)  \cong\mathbb{H}\otimes\mathcal{B}\left(  \mathbb{R}%
^{4}\right)  =M\left(  4,\mathbb{H}\right)
\end{equation}%
\begin{equation}
C\left(  0,6\right)  \cong\mathcal{B}\left(  \mathbb{R}^{8}\right)
\end{equation}%
\begin{equation}
C\left(  7,0\right)  \cong\mathbb{C\otimes H}\otimes\mathcal{B}\left(
\mathbb{R}^{4}\right)
\end{equation}%
\begin{equation}
C\left(  0,7\right)  \cong\mathcal{B}\left(  \mathbb{R}^{8}\right)
\oplus\mathcal{B}\left(  \mathbb{R}^{8}\right)
\end{equation}%
\begin{equation}
C\left(  8,0\right)  \cong C\left(  0,8\right)  \cong\mathcal{B}\left(
\mathbb{R}^{16}\right)
\end{equation}


\subsection{Definite Real Clifford Algebras of Higher Dimension}

All Clifford algebras $C\left(  m,0\right)  ,C\left(  0,m\right)  $ are
determined by $C\left(  n,0\right)  ,C\left(  0,n\right)  ;1\leq n\leq8;0\leq
m<\infty$ by
\begin{equation}
C\left(  n+8,0\right)  \cong C\left(  n,0\right)  \otimes\mathcal{B}\left(
\mathbb{R}^{16}\right)  \subset M\left(  16,C\left(  n,0\right)  \right)
\end{equation}%
\begin{equation}
C\left(  0,n+8\right)  \cong C\left(  0,n\right)  \otimes\mathcal{B}\left(
\mathbb{R}^{16}\right)  \subset M\left(  16,C\left(  0,n\right)  \right)  .
\end{equation}
$\lceil$%
\begin{align}
C\left(  0,15\right)   &  \cong C\left(  0,7\right)  \otimes\mathcal{B}\left(
\mathbb{R}^{16}\right)  \subset M\left(  16,C\left(  0,7\right)  \right) \\
C\left(  0,22\right)   &  \cong C\left(  0,14\right)  \otimes\mathcal{B}%
\left(  \mathbb{R}^{16}\right)  \cong C\left(  0,6\right)  \otimes
\mathcal{B}\left(  \mathbb{R}^{16}\right)  \otimes\mathcal{B}\left(
\mathbb{R}^{16}\right)  \subset M\left(  16\times16,C\left(  0,6\right)
\right)
\end{align}
$\rfloor$

\subsection{Indefinite Real Clifford Algebras}

Clifford algebras of type $C\left(  p,q\right)  $ are descibed by

\begin{enumerate}
\item $p=q$%
\begin{equation}
C\left(  p,p\right)  \cong\mathcal{B}\left(  \mathbb{\wedge R}^{p}\right)
=\mathcal{B}\left(  2^{2p}\right)
\end{equation}


\item $p>q$%
\begin{equation}
C\left(  p,q\right)  \cong C\left(  q,q\right)  \otimes C\left(  s,0\right)
\cong\mathcal{B}\left(  \mathbb{R}^{q}\right)  \otimes C\left(  s,0\right)
\subset M\left(  q,C\left(  s,0\right)  \right)  ,
\end{equation}
where
\begin{equation}
s=p-q
\end{equation}


\item $p<q$%
\begin{equation}
C\left(  p,q\right)  \cong C\left(  p,p\right)  \otimes C\left(  0,r\right)
\cong\mathcal{B}\left(  \mathbb{R}^{p}\right)  \otimes C\left(  0,r\right)
\subset M\left(  p,C\left(  0,r\right)  \right)  .
\end{equation}%
\begin{equation}
r=p-q.
\end{equation}

\end{enumerate}

These relationships plus the modulo $8$ preceding ones characterize\emph{ all
real }Clifford algebras.

$\lceil$%
\begin{align}
2^{4}  &  =16\\
2^{8}  &  =16\times16=256
\end{align}
$\rfloor$

Thus \emph{all }real Clifford algebras are produced from the Clifford algebras
$\operatorname*{Cl}\left(  p,q,\mathbb{R}\right)  ;$ $0\leq p,q\leq8$, whose
structure is given in the following:%

\begin{align}
&  \left(
\begin{array}
[c]{ccccc}%
8 &
\begin{array}
[c]{c}%
\operatorname*{Cl}\left(  8,0,\mathbb{R}\right) \\
M\left(  16,\mathbb{R}\right)
\end{array}
&
\begin{array}
[c]{c}%
\operatorname*{Cl}\left(  8,1,\mathbb{R}\right) \\
M\left(  16,\mathbb{R}\right)  \oplus M\left(  16,\mathbb{R}\right)
\end{array}
&
\begin{array}
[c]{c}%
\operatorname*{Cl}\left(  8,2,\mathbb{R}\right) \\
M\left(  32,\mathbb{R}\right)
\end{array}
&
\begin{array}
[c]{c}%
\operatorname*{Cl}\left(  8,3,\mathbb{R}\right) \\
M\left(  32,\mathbb{C}\right)
\end{array}
\\
7 &
\begin{array}
[c]{c}%
\operatorname*{Cl}\left(  7,0,\mathbb{R}\right) \\
M\left(  8,\mathbb{R}\right)  \oplus M\left(  8,\mathbb{R}\right)
\end{array}
&
\begin{array}
[c]{c}%
\operatorname*{Cl}\left(  7,1,\mathbb{R}\right) \\
M\left(  16,\mathbb{R}\right)
\end{array}
&
\begin{array}
[c]{c}%
\operatorname*{Cl}\left(  7,2,\mathbb{R}\right) \\
M\left(  16,\mathbb{C}\right)
\end{array}
&
\begin{array}
[c]{c}%
\operatorname*{Cl}\left(  7,3,\mathbb{R}\right) \\
M\left(  16,\mathbb{H}\right)
\end{array}
\\
6 &
\begin{array}
[c]{c}%
\operatorname*{Cl}\left(  6,0,\mathbb{R}\right) \\
M\left(  8,\mathbb{R}\right)
\end{array}
&
\begin{array}
[c]{c}%
\operatorname*{Cl}\left(  6,1,\mathbb{R}\right) \\
M\left(  8,\mathbb{C}\right)
\end{array}
&
\begin{array}
[c]{c}%
\operatorname*{Cl}\left(  6,2,\mathbb{R}\right) \\
M\left(  8,\mathbb{H}\right)
\end{array}
&
\begin{array}
[c]{c}%
\operatorname*{Cl}\left(  6,3,\mathbb{R}\right) \\
M\left(  8,\mathbb{H}\right)  \oplus M\left(  8,\mathbb{H}\right)
\end{array}
\\
5 &
\begin{array}
[c]{c}%
\operatorname*{Cl}\left(  5,0,\mathbb{R}\right) \\
M\left(  4,\mathbb{C}\right)
\end{array}
&
\begin{array}
[c]{c}%
\operatorname*{Cl}\left(  5,1,\mathbb{R}\right) \\
M\left(  4,\mathbb{H}\right)
\end{array}
&
\begin{array}
[c]{c}%
\operatorname*{Cl}\left(  5,2,\mathbb{R}\right) \\
M\left(  4,\mathbb{H}\right)  \oplus M\left(  4,\mathbb{H}\right)
\end{array}
&
\begin{array}
[c]{c}%
\operatorname*{Cl}\left(  5,3,\mathbb{R}\right) \\
M\left(  8,\mathbb{H}\right)
\end{array}
\\
4 &
\begin{array}
[c]{c}%
\operatorname*{Cl}\left(  4,0,\mathbb{R}\right) \\
M\left(  2,\mathbb{H}\right)
\end{array}
&
\begin{array}
[c]{c}%
\operatorname*{Cl}\left(  4,1,\mathbb{R}\right) \\
M\left(  2,\mathbb{H}\right)  \oplus M\left(  2,\mathbb{H}\right)
\end{array}
&
\begin{array}
[c]{c}%
\operatorname*{Cl}\left(  4,2,\mathbb{R}\right) \\
M\left(  4,\mathbb{H}\right)
\end{array}
&
\begin{array}
[c]{c}%
\operatorname*{Cl}\left(  4,3,\mathbb{R}\right) \\
M\left(  8,\mathbb{C}\right)
\end{array}
\\
3 &
\begin{array}
[c]{c}%
\operatorname*{Cl}\left(  3,0,\mathbb{R}\right) \\
\mathbb{H\oplus H}%
\end{array}
&
\begin{array}
[c]{c}%
\operatorname*{Cl}\left(  3,1,\mathbb{R}\right) \\
M\left(  2,\mathbb{C}\right)
\end{array}
&
\begin{array}
[c]{c}%
\operatorname*{Cl}\left(  3,2,\mathbb{R}\right) \\
M\left(  4,\mathbb{R}\right)
\end{array}
&
\begin{array}
[c]{c}%
\operatorname*{Cl}\left(  3,3,\mathbb{R}\right) \\
M\left(  4,\mathbb{R}\right)  \oplus M\left(  4,\mathbb{R}\right)
\end{array}
\\
2 &
\begin{array}
[c]{c}%
\operatorname*{Cl}\left(  2,0,\mathbb{R}\right) \\
\mathbb{H}%
\end{array}
&
\begin{array}
[c]{c}%
\operatorname*{Cl}\left(  2,1,\mathbb{R}\right) \\
M\left(  2,\mathbb{C}\right)
\end{array}
&
\begin{array}
[c]{c}%
\operatorname*{Cl}\left(  2,2,\mathbb{R}\right) \\
M\left(  4,\mathbb{R}\right)
\end{array}
&
\begin{array}
[c]{c}%
\operatorname*{Cl}\left(  2,3,\mathbb{R}\right) \\
M\left(  4,\mathbb{R}\right)  \oplus M\left(  4,\mathbb{R}\right)
\end{array}
\\
1 &
\begin{array}
[c]{c}%
\operatorname*{Cl}\left(  1,0,\mathbb{R}\right) \\
\mathbb{C}%
\end{array}
&
\begin{array}
[c]{c}%
\operatorname*{Cl}\left(  1,1,\mathbb{R}\right) \\
M\left(  2,\mathbb{R}\right)
\end{array}
&
\begin{array}
[c]{c}%
\operatorname*{Cl}\left(  1,2,\mathbb{R}\right) \\
M\left(  2,\mathbb{R}\right)  \oplus M\left(  2,\mathbb{R}\right)
\end{array}
&
\begin{array}
[c]{c}%
\operatorname*{Cl}\left(  1,3,\mathbb{R}\right) \\
M\left(  4,\mathbb{R}\right)
\end{array}
\\
0 &
\begin{array}
[c]{c}%
\operatorname*{Cl}\left(  0,0,\mathbb{R}\right) \\
\mathbb{R}%
\end{array}
&
\begin{array}
[c]{c}%
\operatorname*{Cl}\left(  0,1,\mathbb{R}\right) \\
\mathbb{R\oplus R}%
\end{array}
&
\begin{array}
[c]{c}%
\operatorname*{Cl}\left(  0,2,\mathbb{R}\right) \\
M\left(  2,\mathbb{R}\right)
\end{array}
&
\begin{array}
[c]{c}%
\operatorname*{Cl}\left(  0,3,\mathbb{R}\right) \\
M\left(  2,\mathbb{C}\right)
\end{array}
\\
\left.  p\right.  \left/  q\right.  & 0 & 1 & 2 & 3
\end{array}
\right. \\
&  \left.
\begin{array}
[c]{ccccc}%
\begin{array}
[c]{c}%
\operatorname*{Cl}\left(  8,4,\mathbb{R}\right) \\
M\left(  32,\mathbb{H}\right)
\end{array}
&
\begin{array}
[c]{c}%
\operatorname*{Cl}\left(  8,5,\mathbb{R}\right) \\
M\left(  32,\mathbb{H}\right)  \oplus M\left(  32,\mathbb{H}\right)
\end{array}
&
\begin{array}
[c]{c}%
\operatorname*{Cl}\left(  8,6,\mathbb{R}\right) \\
M\left(  64,\mathbb{H}\right)
\end{array}
&
\begin{array}
[c]{c}%
\operatorname*{Cl}\left(  8,7,\mathbb{R}\right) \\
M\left(  128,\mathbb{C}\right)
\end{array}
&
\begin{array}
[c]{c}%
\operatorname*{Cl}\left(  8,8,\mathbb{R}\right) \\
M\left(  256,\mathbb{R}\right)
\end{array}
\\%
\begin{array}
[c]{c}%
\operatorname*{Cl}\left(  7,4,\mathbb{R}\right) \\
M\left(  16,\mathbb{H}\right)  \oplus M\left(  16,\mathbb{H}\right)
\end{array}
&
\begin{array}
[c]{c}%
\operatorname*{Cl}\left(  7,5,\mathbb{R}\right) \\
M\left(  32,\mathbb{H}\right)
\end{array}
&
\begin{array}
[c]{c}%
\operatorname*{Cl}\left(  7,6,\mathbb{R}\right) \\
M\left(  64,\mathbb{C}\right)
\end{array}
&
\begin{array}
[c]{c}%
\operatorname*{Cl}\left(  7,7,\mathbb{R}\right) \\
M\left(  128,\mathbb{R}\right)
\end{array}
&
\begin{array}
[c]{c}%
\operatorname*{Cl}\left(  7,8,\mathbb{R}\right) \\
M\left(  256,\mathbb{R}\right)
\end{array}
\\%
\begin{array}
[c]{c}%
\operatorname*{Cl}\left(  6,4,\mathbb{R}\right) \\
M\left(  16,\mathbb{H}\right)
\end{array}
&
\begin{array}
[c]{c}%
\operatorname*{Cl}\left(  6,5,\mathbb{R}\right) \\
M\left(  32,\mathbb{C}\right)
\end{array}
&
\begin{array}
[c]{c}%
\operatorname*{Cl}\left(  6,6,\mathbb{R}\right) \\
M\left(  64,\mathbb{R}\right)
\end{array}
&
\begin{array}
[c]{c}%
\operatorname*{Cl}\left(  6,7,\mathbb{R}\right) \\
M\left(  64,\mathbb{R}\right)  \oplus M\left(  64,\mathbb{R}\right)
\end{array}
&
\begin{array}
[c]{c}%
\operatorname*{Cl}\left(  6,8,\mathbb{R}\right) \\
M\left(  128,\mathbb{R}\right)
\end{array}
\\%
\begin{array}
[c]{c}%
\operatorname*{Cl}\left(  5,4,\mathbb{R}\right) \\
M\left(  16,\mathbb{C}\right)
\end{array}
&
\begin{array}
[c]{c}%
\operatorname*{Cl}\left(  5,5,\mathbb{R}\right) \\
M\left(  32,\mathbb{R}\right)
\end{array}
&
\begin{array}
[c]{c}%
\operatorname*{Cl}\left(  5,6,\mathbb{R}\right) \\
M\left(  32,\mathbb{R}\right)  \oplus M\left(  32,\mathbb{R}\right)
\end{array}
&
\begin{array}
[c]{c}%
\operatorname*{Cl}\left(  5,7,\mathbb{R}\right) \\
M\left(  64,\mathbb{R}\right)
\end{array}
&
\begin{array}
[c]{c}%
\operatorname*{Cl}\left(  5,8,\mathbb{R}\right) \\
M\left(  64,\mathbb{C}\right)
\end{array}
\\%
\begin{array}
[c]{c}%
\operatorname*{Cl}\left(  4,4,\mathbb{R}\right) \\
M\left(  16,\mathbb{R}\right)
\end{array}
&
\begin{array}
[c]{c}%
\operatorname*{Cl}\left(  4,5,\mathbb{R}\right) \\
M\left(  16,\mathbb{R}\right)  \oplus M\left(  16,\mathbb{R}\right)
\end{array}
&
\begin{array}
[c]{c}%
\operatorname*{Cl}\left(  4,6,\mathbb{R}\right) \\
M\left(  32,\mathbb{R}\right)
\end{array}
&
\begin{array}
[c]{c}%
\operatorname*{Cl}\left(  4,7,\mathbb{R}\right) \\
M\left(  32,\mathbb{C}\right)
\end{array}
&
\begin{array}
[c]{c}%
\operatorname*{Cl}\left(  4,8,\mathbb{R}\right) \\
M\left(  32,\mathbb{H}\right)
\end{array}
\\%
\begin{array}
[c]{c}%
\operatorname*{Cl}\left(  3,4,\mathbb{R}\right) \\
M\left(  8,\mathbb{R}\right)  \oplus M\left(  8,\mathbb{R}\right)
\end{array}
&
\begin{array}
[c]{c}%
\operatorname*{Cl}\left(  3,5,\mathbb{R}\right) \\
M\left(  16,\mathbb{R}\right)
\end{array}
&
\begin{array}
[c]{c}%
\operatorname*{Cl}\left(  3,6,\mathbb{R}\right) \\
M\left(  16,\mathbb{C}\right)
\end{array}
&
\begin{array}
[c]{c}%
\operatorname*{Cl}\left(  3,7,\mathbb{R}\right) \\
M\left(  16,\mathbb{H}\right)
\end{array}
&
\begin{array}
[c]{c}%
\operatorname*{Cl}\left(  3,8,\mathbb{R}\right) \\
M\left(  16,\mathbb{H}\right)  \oplus M\left(  16,\mathbb{H}\right)
\end{array}
\\%
\begin{array}
[c]{c}%
\operatorname*{Cl}\left(  2,4,\mathbb{R}\right) \\
M\left(  8,\mathbb{R}\right)
\end{array}
&
\begin{array}
[c]{c}%
\operatorname*{Cl}\left(  2,5,\mathbb{R}\right) \\
M\left(  8,\mathbb{C}\right)
\end{array}
&
\begin{array}
[c]{c}%
\operatorname*{Cl}\left(  2,6,\mathbb{R}\right) \\
M\left(  8,\mathbb{H}\right)
\end{array}
&
\begin{array}
[c]{c}%
\operatorname*{Cl}\left(  2,7,\mathbb{R}\right) \\
M\left(  8,\mathbb{H}\right)  \oplus M\left(  8,\mathbb{H}\right)
\end{array}
&
\begin{array}
[c]{c}%
\operatorname*{Cl}\left(  2,8,\mathbb{R}\right) \\
M\left(  16,\mathbb{H}\right)
\end{array}
\\%
\begin{array}
[c]{c}%
\operatorname*{Cl}\left(  1,4,\mathbb{R}\right) \\
M\left(  4,\mathbb{C}\right)
\end{array}
&
\begin{array}
[c]{c}%
\operatorname*{Cl}\left(  1,5,\mathbb{R}\right) \\
M\left(  4,\mathbb{H}\right)
\end{array}
&
\begin{array}
[c]{c}%
\operatorname*{Cl}\left(  1,6,\mathbb{R}\right) \\
M\left(  4,\mathbb{H}\right)  \oplus M\left(  4,\mathbb{H}\right)
\end{array}
&
\begin{array}
[c]{c}%
\operatorname*{Cl}\left(  1,7,\mathbb{R}\right) \\
M\left(  8,\mathbb{H}\right)
\end{array}
&
\begin{array}
[c]{c}%
\operatorname*{Cl}\left(  1,8,\mathbb{R}\right) \\
M\left(  16,\mathbb{C}\right)
\end{array}
\\%
\begin{array}
[c]{c}%
\operatorname*{Cl}\left(  0,4,\mathbb{R}\right) \\
M\left(  2,\mathbb{H}\right)
\end{array}
&
\begin{array}
[c]{c}%
\operatorname*{Cl}\left(  0,5,\mathbb{R}\right) \\
M\left(  2,\mathbb{H}\right)  \oplus M\left(  2,\mathbb{H}\right)
\end{array}
&
\begin{array}
[c]{c}%
\operatorname*{Cl}\left(  0,6,\mathbb{R}\right) \\
M\left(  4,\mathbb{H}\right)
\end{array}
&
\begin{array}
[c]{c}%
\operatorname*{Cl}\left(  0,7,\mathbb{R}\right) \\
M\left(  8,\mathbb{C}\right)
\end{array}
&
\begin{array}
[c]{c}%
\operatorname*{Cl}\left(  0,8,\mathbb{R}\right) \\
M\left(  16,\mathbb{R}\right)
\end{array}
\\
4 & 5 & 6 & 7 & 8
\end{array}
\right)
\end{align}


\subsubsection{Example}

Clifford algebra over Minkowski space $p=3,q=1$ is given by
\begin{equation}
C\left(  3,1\right)  =C\left(  2,0\right)  \otimes\mathcal{B}\left(
\mathbb{\wedge R}^{1}\right)  \cong\mathcal{B}\left(  \mathbb{R}^{2}\right)
\otimes\mathcal{B}\left(  \mathbb{R}^{2}\right)  \cong\mathcal{B}\left(
\mathbb{R}^{4}\right)  .
\end{equation}

\end{document}